%\section{FRFP-aware Prompting Patterns}

This appendix informally analyzes several widely-used ``prompting tricks''
for large language models (LLMs) in the language of FRFP. The goal is not to
argue that these prompts eliminate hallucination---they do not, by the
inevitability results of Appendix~B---but to explain \emph{why} they often
improve finite-horizon behavior by acting on the same structural components
that Phase--3 and Phase--4 identify:

\begin{itemize}
  \item the stochastic dynamics of explicit pipelines (the ``Markov kernel''
        on $N \ltimes E$),
  \item the requirement maps $\mathrm{Req}$ that attach strength to
        different speech acts, and
  \item the human user's own tacit process (degradation, attention, HFD).
\end{itemize}

Throughout, we consider a human--LLM pair as a combined FRFP pipeline:

\begin{itemize}
  \item the LLM is an explicit-only component that maps prompts and context
        in $E$ to new artifacts in $E$;
  \item the human user is an epistemic agent $A_H = (T_H,\preceq_H,\delta_H,T^e_H)$
        in the sense of Appendix~C, with representational backflow
        $\mathrm{RB}_H : E \to T_H$;
  \item the joint interaction unfolds as a stochastic FRFP pipeline over
        extended tacit state $T^e_H$ and explicit space $E$ as in Appendix~B.
\end{itemize}

\subsection{FRFP knobs that prompts can touch}

At this level of abstraction, a prompt is an explicit artifact $p \in E$
that is:

\begin{enumerate}
  \item part of the initial explicit state $X_0$ in $N \ltimes E$, and
  \item an input that selects a particular stochastic kernel on explicit
        trajectories (a particular ``policy'') for the LLM component.
\end{enumerate}

Concretely, prompts can affect:

\begin{description}
  \item[(i) Explicit dynamics $M$:] how many steps the LLM takes, which
    tools or subroutines it invokes, whether it produces intermediate
    plans, and so on. In FRFP terms this is a change to the Markov kernel
    on $N \ltimes E$ (Appendix~B.15).

  \item[(ii) Requirements $\mathrm{Req}$:] the kind of speech acts the LLM
    is invited to perform (hypotheses vs.\ assertions, hedged vs.\ strong
    claims), and the normative bar the human tacit process should apply to
    them. Prompts can de facto lower or raise $\mathrm{Req}(e,c)$ by
    changing the intended force of the output.

  \item[(iii) Human tacit dynamics:] the user's own attention, skepticism,
    and willingness to apply HFD and to consult external resources. Many
    prompts are in fact instructions to the \emph{human} about how to
    relate to the model's output, not just instructions to the model.
\end{description}

The five patterns below can be read as specific combinations of (i)--(iii).

\subsection{The Incentive Prompt}

\begin{quote}
\textbf{The Incentive Prompt.} ``I'll tip you \$200 for a perfect solution
to this problem.''
\end{quote}

\paragraph{Mechanism in FRFP terms.}

\begin{enumerate}
  \item \emph{Effect on explicit dynamics $M$.} The textual pattern
    ``I'll pay X if you do Y'' commonly co-occurs, in pretraining data,
    with high-effort, careful responses. Conditioning on such a prompt
    shifts the induced stochastic kernel on $N \ltimes E$ toward trajectories
    with:
    \begin{itemize}
      \item more internal elaboration steps (plans, justifications),
      \item fewer short one-shot answers.
    \end{itemize}
    In the language of Appendix~B.16, this tends to lower some early
    per-step hazards $h_n$ at the cost of slightly longer trajectories.

  \item \emph{Effect on requirements $\mathrm{Req}$.} The phrase ``perfect
    solution'' implicitly raises the bar: the intended speech act is a
    high-Req artifact. Even though the LLM itself has no true tacit state,
    its training aligns this pattern with stronger internal self-checks
    and more cautious phrasing. On the human side, it is natural to attach
    a higher $\mathrm{Req}_H(e,c)$ to any output produced under such
    framing.

  \item \emph{Effect on human tacit dynamics.} The user who writes
    ``I'll tip you \$200'' is also signaling to themselves that the task is
    important. This typically leads to:
    \begin{itemize}
      \item increased attention and willingness to scrutinize the answer,
      \item reduced effective degradation $\delta_H$ over the session.
    \end{itemize}
    In Phase--3 terms, this temporarily slows the drift of $T^e_H$ and
    makes HFD more likely to be applied before acting.
\end{enumerate}

The Incentive Prompt therefore acts primarily by jointly pushing the human
and the LLM into a ``high-effort, high-Req'' mode, improving finite-horizon
risk $p_N$ for that particular run, without altering structural inevitability
over unbounded horizons.

\subsection{The Challenge Prompt}

\begin{quote}
\textbf{The Challenge.} ``I bet you can't solve this correctly. Prove me
wrong.''
\end{quote}

\paragraph{Mechanism in FRFP terms.}

\begin{enumerate}
  \item \emph{Effect on explicit dynamics $M$.} Patterns such as ``prove
    me wrong'' and ``I bet you can't'' are strongly associated in training
    data with argumentative, proof-like discourse. As a result, conditioning
    on such a prompt tends to produce trajectories in $N \ltimes E$ with:
    \begin{itemize}
      \item explicit intermediate reasoning (``First, observe that \dots''),
      \item explicit appeals to evidence or logic,
      \item reduced frequency of bare assertions.
    \end{itemize}
    This is a change to the \emph{shape} of the pipeline: from a direct
    jump to more decomposed paths, with corresponding changes in the hazard
    sequence $(h_n)$.

  \item \emph{Effect on $\mathrm{Req}$.} The words ``correctly'' and
    ``prove'' invite a stronger epistemic commitment: they indicate that
    mere plausibility is insufficient. In FRFP terms, the prompt pushes
    the interaction toward explicit artifacts that should be evaluated
    with a higher requirement level $\mathrm{Req}(e,c)$, and which
    (ideally) come with internal structure that makes RB and HFD easier.

  \item \emph{Effect on human tacit dynamics.} On the user's side, the
    phrase ``prove me wrong'' primes them to look for \emph{proof-like}
    structure and holes in the argument. This acts as a soft instruction
    to apply HFD more aggressively: not only ``does this sound right?''
    but also ``is the reasoning coherent?''. In tacit terms, the user
    temporarily shifts to a stricter criterion within $T_H$ before
    accepting the artifact.
\end{enumerate}

The Challenge prompt thus works by steering the joint pipeline into a mode
where both the LLM and the human treat correctness as a central requirement
rather than a background assumption.

\subsection{The ``Deep Breath'' Prompt}

\begin{quote}
\textbf{The Deep Breath.} ``Take a deep breath and work through this
step by step.''
\end{quote}

\paragraph{Mechanism in FRFP terms.}

\begin{enumerate}
  \item \emph{Effect on explicit dynamics $M$.} This is the clearest
    example of a \emph{pipeline-shaping} prompt. It explicitly instructs
    the LLM to replace a one-step morphism
    \[
      \tau : \text{question} \to \text{answer}
    \]
    by a multi-step trajectory in $N \ltimes E$:
    \[
      \tau_1 : \text{question} \to \text{plan},\quad
      \tau_2 : \text{plan} \to \text{partial result},\ \dots,\ 
      \tau_k : \text{intermediate state} \to \text{final answer}.
    \]
    Under the stochastic dynamics of Appendix~B, this changes both
    the number of steps $N$ and the per-step hazard profile $(h_n)$:
    per-step requirements are typically reduced, while opportunities for
    error detection (by both model and human) are increased.

  \item \emph{Effect on $\mathrm{Req}$.} Breaking a task into substeps
    often lowers the requirement associated with each individual artifact:
    it is easier to check ``is this algebraic manipulation correct?''
    than ``is the whole proof correct?''. The prompt therefore naturally
    redistributes total requirement across the trajectory rather than
    concentrating it in a single, high-Req jump.

  \item \emph{Effect on human tacit dynamics.} Step-by-step reasoning
    aligns better with human cognitive limits: it presents a sequence of
    smaller, more transparent artifacts to $\mathrm{RB}_H$ for evaluation.
    This strengthens the human's ability to apply HFD locally and to
    correct errors before they propagate. In terms of Appendix~B.18, the
    ``Deep Breath'' prompt is an explicit attempt to exploit a more
    favorable time--decomposition trade-off.
\end{enumerate}

Empirically, such prompts are among the most effective at improving
finite-horizon correctness precisely because they operate on the
\emph{structure} of the explicit pipeline rather than only on its tone.

\subsection{The ``Stakes'' Prompt}

\begin{quote}
\textbf{The Stakes.} ``This is very important to my career. You'd better
be sure.''
\end{quote}

\paragraph{Mechanism in FRFP terms.}

\begin{enumerate}
  \item \emph{Effect on explicit dynamics $M$.} Phrases emphasizing high
    stakes and caution are heavily represented in data where the appropriate
    response is:
    \begin{itemize}
      \item hedging,
      \item suggesting external verification, or
      \item narrowing the scope of the answer.
    \end{itemize}
    Conditioning on such a prompt tends to move the induced kernel $M$
    toward:
    \begin{itemize}
      \item more conservative speech acts,
      \item explicit caveats, and
      \item sometimes explicit recommendations to consult a human expert.
    \end{itemize}

  \item \emph{Effect on $\mathrm{Req}$.} This prompt foregrounds the
    requirement function itself: the human is explicitly stating that the
    context $c$ is high-stakes. In a normative FRFP reading, the
    requirement $\mathrm{Req}_H(e,c)$ for any endorsed artifact should
    therefore be substantially higher than in low-stakes chat. The LLM
    approximates this by increasing the proportion of outputs that encode
    uncertainty or suggest further checks.

  \item \emph{Effect on human tacit dynamics.} Perhaps most importantly,
    the prompt serves as a self-reminder for the user: it makes salient
    that their own tacit degradation $\delta_H$ (tiredness, wishful
    thinking) ought to be resisted in this interaction. Users who write
    or read such prompts are more likely to:
    \begin{itemize}
      \item cross-check with other sources,
      \item ask follow-up questions,
      \item avoid taking the output as decisive.
    \end{itemize}
    This is exactly the sort of tacit-side intervention that Appendix~B
    identifies as necessary for any nontrivial mitigation.
\end{enumerate}

\subsection{The Self-Check Prompt}

\begin{quote}
\textbf{The Self-Check.} ``After answering, rate your confidence~0--1.
If below~0.9, try again.''
\end{quote}

\paragraph{Mechanism in FRFP terms.}

\begin{enumerate}
  \item \emph{Effect on explicit dynamics $M$.} This prompt appends two
    additional stages to the explicit pipeline:
    \begin{enumerate}
      \item answer generation: $e \in E$,
      \item self-evaluation: an explicit scalar $\pi(e) \in [0,1]$,
      \item conditional revision: if $\pi(e) < 0.9$, produce $e'$.
    \end{enumerate}
    In $N \ltimes E$ this is a larger composite morphism than a one-shot
    answer. It uses the model's internal heuristics about when it is
    likely to be wrong to steer the dynamics away from low-confidence
    regions of state space.

  \item \emph{Effect on $\mathrm{Req}$.} The scalar $\pi(e)$ functions
    as an explicit \emph{proxy} for how demanding the human should be.
    A low reported confidence is a signal that the artifact should be
    evaluated under a stricter requirement, or even rejected outright.
    This puts a (noisy) estimate of internal plausibility into explicit
    space where the human can use it in their Req\(_H\).

  \item \emph{Effect on human tacit dynamics.} The presence of a reported
    confidence level gives the human a numerical hook for HFD:
    \begin{itemize}
      \item low $\pi(e)$: treat the output as a suggestion, not a belief;
      \item high but not maximal $\pi(e)$: treat with caution in high-risk
            contexts.
    \end{itemize}
    In practice, users often treat confidence scores as prompts to consult
    other tools or experts, which again corresponds to modifying the tacit
    process $T^e_H$ in the way Phase--3 recommends (additional grounding,
    reduced degradation).
\end{enumerate}

\subsection{Limitations and Structural Perspective}

From the FRFP perspective, all five prompt patterns can be summarized as:

\begin{itemize}
  \item reshaping the explicit pipeline in $N \ltimes E$ (more steps,
        different kernels $M$),
  \item altering the effective requirement schedule $\mathrm{Req}(e,c)$
        by changing the intended strength of claims, and
  \item nudging the human tacit process toward more attentive and skeptical
        states in $T^e_H$ for the duration of the interaction.
\end{itemize}

They are therefore valuable \emph{finite-horizon} interventions: for a given
task and risk budget $\varepsilon$, they can materially change the hazard
profile $(h_n)$ and enlarge the safe horizon $H_P(\varepsilon)$ of the joint
pipeline (Appendix~B.16). However, they do not and cannot:

\begin{itemize}
  \item give the LLM genuine grounding beyond the text logs it was trained on,
  \item alter the fundamental explicit--tacit separation (ETS),
  \item or move the combined human--LLM system out of the regime where
        hallucination is inevitable over unbounded horizons (Appendix~B.29
        and Appendix~C.11).
\end{itemize}

In this sense, there is no single ``clever prompt'' that can make hallucination
impossible. What FRFP does provide is a way to classify and compare prompt
design patterns by which structural levers they pull, and to understand why
some patterns---especially those that change the pipeline structure and
human oversight, rather than only the tone of the answer---tend to be more
effective in practice.